\documentclass[12pt]{article}
\usepackage[utf8]{inputenc}
\usepackage[spanish]{babel}
\usepackage{amsmath}
\usepackage{amssymb}
\usepackage{braket} % Paquete ideal para la notación de Dirac
\usepackage{geometry}
\geometry{a4paper, margin=2.5cm}

\title{Sobre los primeros pasos del caminante de una moneda de 4D}
\author{Saúl Nájera Allara}
\date{}

\begin{document}

\maketitle

Sea $\mathcal{H}_p$ el espacio de hilbert de la posición y $\mathcal{H}_c$ el espacio de hilbert de la moneda. El estado del caminante 
pertenece al espacio:
\begin{equation*}
\mathcal{H}=\mathcal{H}_p \otimes \mathcal{H}_c
\end{equation*}
Considerando que $dim\, \mathcal{H}_c=4$ y sea $\{c_i\}_{i=1,\cdots,4}$ una base ortonormal de $\mathcal{H}_c$.
Podemos interpretar que $\ket{c_1}$ denota la dirección izquierda, $\ket{c_2}$ arriba, $\ket{c_3}$ derecha, $\ket{c_4}$ abajo.

Dado $\ket{c}\in \mathcal{H}_c$, se sigue que:
\begin{equation*}
    \ket{c}=\sum_{m=1}^{4}\lambda_m \ket{c_m}.
\end{equation*}
Pidiendo que $\bra{c}\ket{c}=1$, se obtiene:
\begin{equation*}
    \sum_{m=1}^{4}|\lambda_n|^2=1.
\end{equation*}
Un estado localizado genérico es de la forma:
\begin{equation}
    \ket{\psi}=\sum_{m=1}^{4}\ket{x,y}\otimes\ket{\lambda_m c_m}.
    \label{eq:eq1}
\end{equation}
Si $\Omega$ denota el interior del estadio, entonces se define la acción del operador shift condicional S sobre un estado de la forma dada en \eqref{eq:eq1} como:
\begin{equation}
   S\ket{x,y}\otimes \ket{c_m}=
   \begin{cases}
        \ket{x+\delta_m^{(x)},y+\delta_m^{(y)}}\otimes \ket{c_m} & \text{si $(x,y)\in \Omega$}.
        \\
        \ket{x,y}\otimes \ket{\bar{c}_m} & \text{si $(x,y)\notin \Omega$}.
    \end{cases}
\end{equation}
Donde $\delta_m^{(x)}$, $\delta_m^{(x)}$, denotan desplazamientos unitarios (se desplazan a alguno de los vértices más cercanos) que depende de una regla fija asociada los $c_m$
y $\bar{c_m}$ denota invertir la dirección de $c_m$, es decir, izquierda pasa a derecha, arriba pasa a abajo, etc.
En particular, se considera la regla de asignación según:
\begin{equation*}
\begin{cases}
    \text{Si $\ket{c_m}=\ket{c_1}$} & \text{, entonces }\delta_1^{(x)}=-1\,,\delta_1^{(y)}=0. \\
     \text{Si $\ket{c_m}=\ket{c_2}$} & \text{, entonces }\delta_2^{(x)}=0\,,\delta_2^{(y)}=1. \\
      \text{Si $\ket{c_m}=\ket{c_3}$} & \text{, entonces }\delta_3^{(x)}=1\,,\delta_3^{(y)}=0. \\
       \text{Si $\ket{c_m}=\ket{c_4}$} & \text{, entonces }\delta_4^{(x)}=0\,,\delta_4^{(y)}=-1. \\
\end{cases}
\end{equation*}
Defínase el operador moneda según:
\begin{equation}
   C(\theta)= 1_p\otimes M(\theta)
\end{equation}
Donde $M(\theta)$ viene dada según:
\begin{equation*}
M(\theta) = 
\begin{pmatrix}
-\sin^2\theta & \cos^2\theta & \cos\theta\sin\theta & \cos\theta\sin\theta \\
\cos^2\theta & -\sin^2\theta & \cos\theta\sin\theta & \cos\theta\sin\theta \\
\cos\theta\sin\theta & \cos\theta\sin\theta & -\cos^2\theta & \sin^2\theta \\
\cos\theta\sin\theta & \cos\theta\sin\theta & \sin^2\theta & -\cos^2\theta
\end{pmatrix}
\end{equation*}
En particular, para $\theta=\frac{\pi}{4}$, se tiene:
\begin{equation*}
M(\frac{\pi}{4}) = 
\begin{pmatrix}
-1/2 & 1/2 & 1/2 & 1/2 \\
1/2 & -1/2 & 1/2 & 1/2 \\
1/2 & 1/2 & -1/2 & 1/2 \\
1/2 & 1/2 & 1/2 & -1/2
\end{pmatrix}
\end{equation*}
$M(\frac{\pi}{4})$ es unitaria.

Si consideramos que los $\ket{c_i}$ son los vectores canónicos $\ket{e_i}$, entonces se tiene para $\theta=\frac{\pi}{4}$ que:
\begin{align*}
\begin{cases}
M(\frac{\pi}{4})\ket{c_1}=-\frac{1}{2}\ket{c_1}+\frac{1}{2}\ket{c_2}+\frac{1}{2}\ket{c_3}+\frac{1}{2}\ket{c_4}
\\
M(\frac{\pi}{4})\ket{c_2}=\frac{1}{2}\ket{c_1}-\frac{1}{2}\ket{c_2}+\frac{1}{2}\ket{c_3}+\frac{1}{2}\ket{c_4}
\\
M(\frac{\pi}{4})\ket{c_3}=\frac{1}{2}\ket{c_1}+\frac{1}{2}\ket{c_2}-\frac{1}{2}\ket{c_3}+\frac{1}{2}\ket{c_4}
\\
M(\frac{\pi}{4})\ket{c_4}=\frac{1}{2}\ket{c_1}+\frac{1}{2}\ket{c_2}+\frac{1}{2}\ket{c_3}-\frac{1}{2}\ket{c_4}
\end{cases}
\end{align*}
Vemos que
\begin{equation*}
M(\theta)\ket{c_i}=\sum_{k=1}^{4}\alpha_k \ket{c_k}\, , \alpha_k \neq 0
\end{equation*}
Note que al ser unitaria, se preserva la norma y por tanto $\alpha_i$ son amplitudes de probabilidad.
El operador evolución se define según:
\begin{equation}
U=S\cdot C(\theta).
\end{equation}
Consideremos el primer paso de un caminante de la forma:
\begin{equation}
    \ket{\psi_0}=\ket{x,y}\otimes\ket{c_1}
\end{equation}
Entonces:
\begin{align*}
\ket{\psi_1}= S\cdot C(\theta)\ket{\psi_0}=S\cdot C(\theta)\ket{x,y}\otimes\ket{c_1}
\\
=S\ket{x,y}\otimes M(\theta)\ket{c_1}
\\
=S\ket{x,y}\otimes \sum_{k=1}^{4}\alpha_k \ket{c_k}
\\
\ket{\psi_1}=\sum_{k=1}^{4}S\ket{x,y}\otimes \alpha_k \ket{c_k}
\end{align*}
Aplicando las reglas de asignación establecidas, se tendría que:
\begin{equation}
    \ket{\psi_1}=\ket{x-1,y}\otimes \alpha_1 \ket{c_1}+\ket{x,y+1}\otimes \alpha_2 \ket{c_2}+\ket{x+1,y}\otimes \alpha_3 \ket{c_3}
    +\ket{x,y-1}\otimes \alpha_4 \ket{c_4}
\end{equation}
Inmediantamente podemos concluir que la evolución a un primer paso de cualquier estado localizado de la forma \eqref{eq:eq1} se obtienen cuatro estados localizados en las casillas inmediatas a la posición inicial
y para cada estado se tiene como estado de moneda la dirección en la que se movió con un peso $\alpha_k$ de que suceda.
Las $\alpha_i$ dependen directamente de $M(\theta)$ y para el caso de $\theta=\frac{\pi}{4}$ se para cada dirección posible, igual probabilidad.

Para estados de la forma 
\begin{equation}
    \ket{\phi}=\ket{x,y}\otimes \lambda \ket{c_i}
\end{equation}
Considere 
\begin{equation*}
    \ket{\phi}=(1_p\otimes \lambda)\ket{x,y}\otimes \ket{c_i}
\end{equation*}
Y se verifica que: \[[U,1_p\otimes \lambda]=0\]
Y aplicando el mismo análisis que antes, se sigue que, para el estado
\[\ket{\phi_0}=\ket{x,y}\otimes \lambda \ket{c_1}\]
Se tiene la evolución
\begin{align*}
   \ket{\phi_1}= U\ket{\phi_0}=(1_p\otimes \lambda )U\ket{x,y}\otimes \ket{c_1}
   \\
   \\
  \ket{\phi_1} =(1_p\otimes \lambda )(\ket{x-1,y}\otimes \alpha_1 \ket{c_1}+\ket{x,y+1}\otimes \alpha_2 \ket{c_2}
   \\
   +\ket{x+1,y}\otimes \alpha_3 \ket{c_3}+\ket{x,y-1}\otimes \alpha_4 \ket{c_4})
   \\\\
    \ket{\phi_1} =\ket{x-1,y}\otimes \lambda\alpha_1 \ket{c_1}+\ket{x,y+1}\otimes \lambda\alpha_2 \ket{c_2}
   \\
   +\ket{x+1,y}\otimes \lambda\alpha_3 \ket{c_3}+\ket{x,y-1}\otimes \lambda\alpha_4 \ket{c_4}
\end{align*}

\end{document}